\chapter{Methodology}
\label{ch:methodology}

\section{Research Design}
\label{sec:design}
The study uses a quantitative, experiment-driven design. All four model families are framed as the
same task, \emph{bounded ordinal regression}, and trained and evaluated through one shared,
Khmer-aware pipeline, so that differences in results come from the models rather than from
inconsistent preprocessing or evaluation. Each model predicts a continuous normalised score
$\hat{y}\in[0,1]$, which is mapped to an integer grade for reporting. Every model is then judged on
two axes: standard agreement metrics \emph{and} the faithfulness of its explanations. The overall
pipeline is shown in Figure~\ref{fig:method}.

\begin{figure}[t]
\centering
\begin{tikzpicture}[
  font=\footnotesize, >={Latex[length=2mm]},
  wide/.style={draw, rounded corners, align=center, inner sep=3pt, fill=blue!4, text width=12cm},
  half/.style={draw, rounded corners, align=center, inner sep=3pt, fill=blue!4, text width=6.0cm},
  pill/.style={draw, rounded corners, align=center, inner sep=2pt, fill=blue!8,
               text width=2.7cm, minimum height=10mm},
  io/.style={draw, rounded corners, align=center, inner sep=3pt, fill=gray!10, text width=12.8cm},
  arr/.style={->, thick}]
  \node[io]   (data)  at (0,0)      {\textbf{Khmer ASAG dataset}, 1{,}184 answers, 41 questions, 4 subjects \;[\,Question, Reference, Student answer, Teacher score, Max score\,]};
  \node[wide] (pre)   at (0,-1.5)   {\textbf{Preprocessing $\Phi$ (Khmer-aware)}: NFC $\to$ strip invisibles $\to$ KCC reorder $\to$ strip punctuation $\to$ (khmernltk segmentation)};
  \node[wide] (split) at (0,-2.9)   {\textbf{Split 70/15/15} (seed 42): random (stratified) $\;|\;$ question-held-out (GroupShuffleSplit by QuestionID)};
  \node[pill] (cl)  at (-5.4,-4.7) {\textbf{Classical}\\ TF-IDF + RBF-SVR};
  \node[pill] (rnn) at (-1.8,-4.7) {\textbf{RNN}\\ BiLSTM + Attn};
  \node[pill] (tr)  at (1.8,-4.7)  {\textbf{Transformer}\\ mBERT/XLM-R/GTE};
  \node[pill] (llm) at (5.4,-4.7)  {\textbf{LLM}\\ Qwen (QLoRA)};
  \node[wide] (score) at (0,-6.5) {\textbf{Score $\to$ grade}: round$(4\hat{y})$, $\hat{y}\in[0,1]$ \;(+ optional threshold calibration, selected on validation)};
  \node[half] (eval) at (-3.4,-8.4) {\textbf{Evaluation}\\ QWK, Accuracy, macro-F1, Cohen $\kappa$;\\ exact / adjacent ($\pm1$) agreement};
  \node[half] (xai)  at (3.4,-8.4)  {\textbf{Explainability}\\ word attribution: Leave-One-Out (LOO)\\ $\to$ ERASER faithfulness + plausibility};
  \node[io] (proto) at (0,-10.2) {\textbf{Live web prototype}: teacher enters Question + Reference + Answer $\to$ choose model $\to$ \textbf{score + explanation + feedback}};
  \draw[arr] (data)--(pre);
  \draw[arr] (pre)--(split);
  \foreach \p in {cl,rnn,tr,llm} {\draw[arr] (split)--(\p);}
  \foreach \p in {cl,rnn,tr,llm} {\draw[arr] (\p)--(score);}
  \draw[arr] (score)--(eval);
  \draw[arr] (score)--(xai);
  \draw[arr] (eval)--(proto);
  \draw[arr] (xai)--(proto);
\end{tikzpicture}
\caption{Methodology overview. Four interchangeable model pillars are trained as bounded ordinal
regression on one Khmer-aware pipeline, then evaluated on standard metrics \emph{and} for
explanation faithfulness; the system is realised as a live teacher-facing web prototype.}
\label{fig:method}
\end{figure}

\section{Dataset}
\label{sec:dataset}
The corpus consists of \textbf{1{,}184} short answers (after dropping one incomplete row) collected
from two schools, eight classes, and 203 students, covering \textbf{41} distinct questions across
four secondary-school subjects (Biology, History, Geography, and Earth Science). Each record contains
the question, a reference answer, the student's answer, the teacher's integer score, and the
question's maximum score. A single trained teacher graded every answer. The per-question maximum
score varies across $\{5,6,7,8,10,12,15,20\}$.

To place all questions on a common scale, each answer is given a \emph{normalised score}
\begin{equation}
y \;=\; \frac{\text{Student Score}}{\text{Max Score}} \;\in\; [0,1],
\end{equation}
and an ordinal five-class \emph{grade label}
\begin{equation}
\ell \;=\; \mathrm{clip}\big(\mathrm{round}(4y),\,0,\,4\big)\;\in\;\{0,1,2,3,4\}.
\end{equation}
The label distribution is heavily skewed, with roughly 42\% of answers receiving full credit.

Three curated \emph{variants} of the corpus are used so that the effect of data quality can be
measured (Table~\ref{tab:variants}): the full set; a version with a noisy ``10C Biology'' subset
removed; and a version that additionally drops 14 answers whose grade label is $0$ (about 1.2\% of
the data, effectively untrainable label noise). All three are reported, so the effect of these
choices is visible rather than hidden.

\begin{table}[t]
\centering
\caption{The three curated dataset variants.}
\label{tab:variants}
\begin{tabular}{@{}llr@{}}
\toprule
Variant & Definition & Rows \\
\midrule
\texttt{full}        & Original corpus                                  & 1{,}184 \\
\texttt{no10c}       & Drop the noisy ``10C Biology'' subset            & 909 \\
\texttt{no10c\_no0}  & Also drop the 14 grade-$0$ outliers              & 895 \\
\bottomrule
\end{tabular}
\end{table}

Unless stated otherwise, the data is split 70/15/15 into training, validation, and test sets with a
fixed random seed (42). Two split modes are used: a \emph{stratified random} split over rows (the
default, comparable to prior work) and a stricter \emph{question-held-out} split (a
\texttt{GroupShuffleSplit} by question identifier) in which no question is shared between train,
validation, and test (\S\ref{sec:metrics}).

\section{Preprocessing}
\label{sec:preprocessing}
Khmer text needs careful normalisation. The pipeline applies, in order: Unicode NFC normalisation;
removal of \emph{invisible} characters (zero-width spaces, format and control characters, and bullet
markers such as \km{•} \km{◦}) that are genuine noise surviving punctuation stripping; \emph{KCC}
(Khmer Character Cluster) reordering, which canonicalises the order of base consonants, subscript
consonants (\km{ជើង}), vowels, and signs within each cluster; punctuation stripping; and finally word
segmentation with \texttt{khmer-nltk}. Digits (both Arabic and Khmer numerals) and letters are kept
as answer content. Three preprocessing modes are exposed, \texttt{raw} (normalisation only),
\texttt{clean} (through punctuation stripping), and \texttt{segment} (adds word segmentation), and
treated as an axis of the experiment grid. The released headline results were produced with the
pre-refinement cleaning; an ablation (Chapter~\ref{ch:results}) shows that the later
invisible-stripping refinement changes QWK by less than 0.01.

\section{The Four Model Pillars}
\label{sec:pillars}
All models consume the same answer/reference text pair (optionally prefixed with the question) and
output $\hat{y}\in[0,1]$, trained to minimise mean-squared error against $y$.

\subsection{Classical: TF-IDF features with RBF-SVR}
The classical pillar represents each text with character-level TF-IDF features (using
\texttt{char\_wb} $n$-grams of length 2--4, capped at 15{,}000 features), which suit Khmer's
space-free script. For an answer/reference pair with feature vectors $\mathbf{a}$ and $\mathbf{b}$,
the model forms the interaction vector
\begin{equation}
\big[\,\mathbf{a};\ \mathbf{b};\ |\mathbf{a}-\mathbf{b}|;\ \mathbf{a}\odot\mathbf{b};\
\cos(\mathbf{a},\mathbf{b})\,\big],
\end{equation}
and feeds it to a support-vector regressor with an RBF kernel. A cosine-similarity baseline is also
included. This pillar trains in seconds on a CPU.

\subsection{RNN: character BiLSTM with attention}
The recurrent pillar embeds each text at the character level and encodes it with a bidirectional LSTM.
An additive attention layer pools the per-position hidden states into a single vector; concretely, for
hidden states $\mathbf{h}_t$ the attention weights are $\alpha_t=\mathrm{softmax}(\mathbf{w}^\top
\tanh(W\mathbf{h}_t))$ and the encoding is $\sum_t \alpha_t \mathbf{h}_t$. The answer and reference
are encoded independently and combined through the same four-way interaction
$[\mathbf{e}_a;\mathbf{e}_b;|\mathbf{e}_a-\mathbf{e}_b|;\mathbf{e}_a\odot\mathbf{e}_b]$ before a small
multilayer-perceptron head with a sigmoid output. The attention mechanism is part of the encoding architecture; it is not used as an explanation.

\subsection{Transformer: multilingual BERT-family encoders}
The Transformer pillar fine-tunes multilingual encoders, mBERT, XLM-R, and
GTE-multilingual~\citep{devlin2019bert,conneau2020xlmr}, in two configurations: a \emph{dual}
encoder that embeds answer and reference separately and combines them, and a \emph{cross} encoder
that feeds both into a single sequence and lets self-attention relate them. The bottom six layers are
frozen; inputs are truncated to 256 tokens. (GTE-multilingual ships a numerically unstable rotary
position cache; the implementation rebuilds it in full precision on load to avoid NaN outputs.)

\subsection{LLM: Qwen 3.5 4B fine-tuned with QLoRA}
The largest pillar fine-tunes the Qwen 3.5 4B large language model with QLoRA (quantised low-rank adaptation),
which trains a small number of low-rank adapter weights on top of a frozen, 4-bit-quantised base
model. The model is prompted with the question, reference, and student answer, and trained to emit the
score; it is scored, and LOO word attribution is applied identically to all other families.

\section{From Score to Grade: the Max-Score Feature and Calibration}
\label{sec:score-to-grade}
The continuous prediction is mapped to an integer grade by rounding $4\hat{y}$ (for the five-class
label) or $\hat{y}\times\text{Max Score}$ (for the raw integer score). Two low-cost levers are
studied (RQ4). First, because a prediction of ``0.5'' means 5/10 on one question but 10/20 on another,
the question's maximum score is supplied to the model as an extra normalised input feature
($\text{Max Score}/20$); this gives consistent gains on the deployment metrics. Second, \emph{post-hoc
threshold calibration} fits the four cut-points that map $\hat{y}$ to the five classes by maximising
QWK \emph{on the validation set} (coordinate descent). Calibration is treated as an \emph{ablation},
not part of the headline: as Chapter~\ref{ch:results} shows, the same validation-selected rule helps
the classical model but \emph{hurts} the BiLSTM on the test set, so it is a fragile, model-dependent
lever and all headline numbers are reported \emph{uncalibrated}.

\section{Explainability Protocol}
\label{sec:xai}
Explanations are produced and \emph{evaluated} at the level of Khmer \emph{word units} (from the
segmenter) on the same test answers, so that families are directly comparable. The reference side is
held fixed; only the answer is attributed, because that is what a teacher grades.

\paragraph{Word attribution (text highlighting).} The explainability contribution is delivered as
\emph{word attribution}: the system highlights exactly which words in the student's answer drove the
predicted grade. The model-agnostic engine is \emph{occlusion} (leave-one-out, LOO), the standard ASAG
attribution~\citep{pinto2025attribution}; the importance of word $i$ is
\begin{equation}
\mathrm{imp}(i) \;=\; s(\text{answer}) \;-\; s(\text{answer without word } i),
\end{equation}
where $s(\cdot)$ is the model's predicted score and a positive value means the word supported the
grade. LOO occlusion is the \emph{sole} explanation method because it works for every pillar, including
the non-differentiable classical SVR, requires no gradients or internal model access, is
deterministic, and is directly aligned with the ERASER faithfulness metrics used to evaluate it.

\paragraph{Two evaluation axes.} Following \citet{jacovi2020faithfulness}, explanations are judged on
two axes. \emph{Faithfulness} asks whether the explanation reflects the model's actual reasoning,
measured by the ERASER metrics~\citep{deyoung2020eraser}: \emph{comprehensiveness} (removing the
top-$k$ most important words should drop the score, higher is better) and \emph{sufficiency} (keeping
only the top-$k$ words should preserve the score, lower is better), each compared against a
\emph{random-removal} baseline and averaged over $k\in\{10,20,30,40,50\}\%$ as the AOPC summary
\citep{samek2017aopc}. A positive \emph{faithfulness gap} (comprehensiveness minus the random
baseline) means the explanation beats removing random words. \emph{Plausibility} asks whether the
explanation looks reasonable, approximated by \emph{reference overlap}: the fraction of the top
important answer words that also appear in the teacher's reference answer. Plausibility is an
automatic proxy, not a human-rationale study. 
\section{Evaluation Methodology}
\label{sec:metrics}
\paragraph{Agreement metrics.} The primary metric is \emph{Quadratic Weighted Kappa} (QWK), the
field standard for ordinal scoring~\citep{cohen1968qwk}. For confusion matrix $O$, expected matrix
$E$, and quadratic weights $w_{ij}=(i-j)^2/(K-1)^2$,
\begin{equation}
\mathrm{QWK} \;=\; 1 - \frac{\sum_{i,j} w_{ij}\,O_{ij}}{\sum_{i,j} w_{ij}\,E_{ij}}.
\end{equation}
For direct comparison with the Arabic anchor~\citep{alaoui2024}, \emph{unweighted Cohen's $\kappa$} is
also reported, along with \emph{accuracy} and \emph{macro-averaged F1}. For deployment, following
\citet{williamson2012framework}, \emph{exact agreement} (predicted integer score equals the teacher's)
and \emph{adjacent agreement} (within one point) are reported.

\paragraph{Honest reporting.} Every metric is computed from each run's saved per-answer predictions,
so all results can be reproduced without re-training, and every cell is reported on both the train and
test splits across all three dataset variants, which exposes any overfitting and lets us attribute
gains to data quality versus algorithm rather than to a single favourable run.

\paragraph{Question-held-out evaluation.} To quantify question leakage (RQ6), the classical champion
is re-evaluated under the question-held-out split (no question shared between train/validation/test),
across five seeds. With only 41 questions, each unseen-question test set contains roughly six
questions, so this estimate has high variance, which is reported.

\section{Ethical Considerations}
\label{sec:ethics}
The corpus contains real student work by minors. Only pseudonymous identifiers are stored (no
names), and the data is used solely for research into a teacher-assist tool, not an autonomous
grader that replaces teachers. Because a single teacher graded all answers, the system reflects that
grader's judgements and may carry their biases; results are reported as agreement \emph{with this
grader}, with reporting across three dataset variants as partial mitigation. The full data-governance statement is
included in Appendix~\ref{app:ethics}.
